\documentclass[14pt,a4paper]{extarticle}
\usepackage[utf8]{inputenc}
\usepackage[T2A]{fontenc}
\usepackage[russian]{babel}
\usepackage{amsmath}
\usepackage{geometry}
\usepackage{hyperref}
\geometry{margin=2cm}

\title{Лабораторная работа №5\\Интерполяция функции}
\author{Шубин Е.\,В., P3209}
\date{2026}

\begin{document}
\maketitle

\section*{Цель работы}
Решить задачу интерполяции табличной функции и сравнить результаты, полученные разными интерполяционными многочленами.

\section*{Порядок выполнения}
\begin{enumerate}
    \item Реализована программа с GUI на Python.
    \item Добавлены три способа ввода данных: вручную, из файла, генерация по функции.
    \item Реализованы методы Лагранжа, Ньютона (разделенные и конечные разности), Гаусса.
    \item Выводятся таблицы конечных и разделенных разностей.
    \item Выполняется вычисление значения функции в точке $x^*$.
    \item Строится график узлов и интерполяционных полиномов.
\end{enumerate}

\section*{Рабочие формулы}
\textbf{Лагранж:}
\[
L_n(x)=\sum_{i=0}^{n}y_i\prod_{j=0,j\ne i}^{n}\frac{x-x_j}{x_i-x_j}.
\]

\textbf{Ньютон (разделенные разности):}
\[
N_n(x)=f[x_0]+f[x_0,x_1](x-x_0)+\dots+f[x_0,\dots,x_n]\prod_{k=0}^{n-1}(x-x_k).
\]

\textbf{Ньютон (конечные разности, вперед):}
\[
N_n(x)=y_0+t\Delta y_0+\frac{t(t-1)}{2!}\Delta^2y_0+\dots,\quad t=\frac{x-x_0}{h}.
\]

\textbf{Гаусс (первая формула):}
\[
P_n(x)=y_0+t\Delta y_0+\frac{t(t-1)}{2!}\Delta^2y_{-1}
+\frac{(t+1)t(t-1)}{3!}\Delta^3y_{-1}+\dots
\]

\section*{Вычислительная часть (вариант 18)}
Табличные данные:
\[
\begin{array}{c|ccccccc}
x & 1.10 & 1.25 & 1.40 & 1.55 & 1.70 & 1.85 & 2.00 \\
\hline
y & 0.2234 & 1.2438 & 2.2644 & 3.2984 & 4.3222 & 5.3516 & 6.3867
\end{array}
\]

Точки интерполяции: $X_1=1.875$, $X_2=1.575$.

В программе для этих точек получаются значения всеми реализованными методами; результаты сравниваются по близости.

\section*{Листинг программы}
Исходный код: \url{https://github.com/<your-username>/Itmo_labs/tree/main/4%20sem%20ComputationalMath/Lab5}

\section*{Результаты программы}
Программа корректно обрабатывает:
\begin{itemize}
    \item табличные данные варианта 18;
    \item тестовые файлы из директории \texttt{test\_data};
    \item точки, сгенерированные по функциям \texttt{sin(x)}, \texttt{cos(x)}, \texttt{exp(x)}.
\end{itemize}

\section*{Вывод}
В ходе работы реализованы и протестированы основные методы интерполяции. Для равноотстоящих узлов удобно использовать конечные разности и формулы Ньютона/Гаусса, а формула Лагранжа универсальна и применима к произвольной сетке.

\end{document}
